% =============================================================
% 子文件模板：题目描述部分
% 用法：在主文件中使用 \input{filename.tex} 导入
% =============================================================

\begin{center}
    % 建议格式：序号. 中文题目名（英文题目名）
    \section{A. 异域旅行 (travel)}
\end{center}

\subsection{题目描述}

% 在此处填写题目背景和详细描述
% 行内公式请使用 $...$，例如 $a+b$
小 A 有 $n$ 只小狗，第 $i$ 只小狗的战斗力为 $a_i$。一天，他们掉到了一个异域之中，于是第 $i$ 只小狗的战斗力将变为 $b_i = a_i \oplus k$。

我定义第 $i$ 只小狗的进步值为 $\max(0,b_i-a_i)$，请你输出所有小狗进步值总和，即为 $\sum_{i=1}^{n} \max(0,b_i-a_i)$。

注：$\oplus$ 指按位异或。

\subsection{输入格式}

% 描述输入数据的格式，例如：第一行包含一个整数 $n$。
第一行，两个正整数，$n$ 和 $k$，含义如题面所示。

第二行 $n$ 个正整数，第 $i$ 个数代表 $a_i$。

\subsection{输出格式}

% 描述输出数据的格式
一个数，代表进步值总和。

\subsection{样例输入1}

% 将样例输入直接粘贴在下方，无需额外格式
\begin{lstlisting}
5 3
0 7 8 4 2
\end{lstlisting}

\subsection{样例输出1}

% 将样例输出直接粘贴在下方
\begin{lstlisting}
9
\end{lstlisting}

\subsection{样例输入2}

% 将样例输入直接粘贴在下方，无需额外格式
\begin{lstlisting}
6 4
1 15 8 9 2 14
\end{lstlisting}

\subsection{样例输出2}

% 将样例输出直接粘贴在下方
\begin{lstlisting}
16
\end{lstlisting}

% 如果有更多样例，可以复制上方的 subsection 块
% \subsection{样例输入2}...

\subsection{数据范围和约定}

% 描述数据规模和特殊限制
% 示例：对于 $100\%$ 的数据，$1 \leq n \leq 100$。
对于 $20\%$ 的数据，$n \le 1000$。

对于另外 $10\%$ 的数据，$k = 0$。

对于 $100\%$ 的数据，$1 \le n \le 10^5,1 \le a_i \le 10^6,1 \le k \le 10^6$。
